
\section{Research Design \& Paradigm}
This study is structured under a \textbf{positivist, quantitative-dominant research design}, employing empirical geospatial observations, multi-spectral band math, and unsupervised machine learning algorithms to evaluate orbital data as a proxy for real-world economic activity. The research paradigm is grounded in the physical sciences: satellite sensors measure raw solar reflectance and microwave backscatter, which are objective physical properties of the Earth's surface. 

These physical observations are processed through a structured Information Engineering pipeline, transforming unstructured raster bands into structured anomaly events. While the core processing and modeling are quantitative, a qualitative ground-truthing layer is integrated to cross-reference detected anomalies with news reports, corporate press releases, and official mining registries, establishing a robust mixed-method validation framework.

---

\section{Data Ingestion \& API Architectures}
The automated ingestion pipeline queries three primary geospatial data portals and one socioeconomic database, operating under rate-limited, decoupled connection parameters:

\subsection{Google Earth Engine (GEE) Ingestion}
The pipeline interfaces with GEE utilizing the Python \texttt{earthengine-api} library. GEE serves as the primary engine for processing large-scale spatial collections. Ingestion parameters are configured dynamically:
\begin{itemize}
  \item \textbf{Coordinate Bounds}: Bounding boxes are defined in WGS 84 (EPSG:4326) coordinates.
  \item \textbf{Cloud Filtering}: Collections are filtered to keep only scenes with cloud cover below 20\% using the \texttt{CLOUDY\_PIXEL\_PERCENTAGE} attribute for Sentinel-2, or \texttt{CLOUD\_COVER} for Landsat-8/9.
  \item \textbf{Temporal Filtering}: Image collections are filtered using start and end date boundaries to establish distinct baseline and anomaly comparison epochs.

\end{itemize}
\subsection{Copernicus Data Space Ecosystem (CDSE) Ingestion}
CDSE OData and Sentinel Hub APIs are queried programmatically to ingest high-resolution Sentinel-1 SAR and Sentinel-2 optical bands.
\begin{enumerate}
  \item \textbf{OAuth Authentication}: The system requests a short-lived bearer token using client credentials.
  \item \textbf{Process API Payload}: A POST request containing the bounding box, output resolution (10m), target bands, and custom Evalscript (which computes indices and cloud masks on the server side) is transmitted.
  \item \textbf{Raster Acquisition}: The CDSE backend returns a single multi-band GeoTIFF or scaled NumPy array, minimizing local bandwidth.

\end{enumerate}
\subsection{Socioeconomic API Integration}
The World Bank API is queried to fetch regional annual GDP metrics. The REST query is structured as:
\texttt{http://api.worldbank.org/v2/country/{country}/indicator/NY.GDP.MKTP.KD.ZG?format=json}
The system extracts the tabular data and matches the temporal records with satellite night-time lights anomalies.

---

\section{Spectral Index Formulations}
To convert raw digital numbers (DN) or top-of-atmosphere (TOA) reflectance values into indicators of land-use change, the pipeline implements four mathematical index formulations:

\subsection{Normalized Difference Vegetation Index (NDVI)}
NDVI quantifies vegetation density and health by contrasting the high absorption of chlorophyll in the red visible band with the high reflectance of leaf mesophyll in the near-infrared (NIR) band:
$$\text{NDVI} = \frac{\rho\_{\text{NIR}} - \rho\_{\text{Red}}}{\rho\_{\text{NIR}} + \rho\_{\text{Red}}}$$
\begin{itemize}
  \item \textbf{Sentinel-2 Band Map}: $\frac{B8 - B4}{B8 + B4}$
  \item \textbf{Landsat-8/9 Band Map}: $\frac{B5 - B4}{B5 + B4}$
  \item \textbf{Reflectance Range}: $[-1.0, 1.0]$. Dense forest canopies exhibit values between $0.6$ and $0.9$. Land clearing, open-pit mining, or infrastructure construction drops NDVI values close to $0.0$ or into negative ranges.

\end{itemize}
\subsection{Normalized Difference Water Index (NDWI)}
NDWI, originally formulated by McFeeters (1996), maps open water surfaces and monitors changes in moisture levels by utilizing the green band and the NIR band:
$$\text{NDWI} = \frac{\rho\_{\text{Green}} - \rho\_{\text{NIR}}}{\rho\_{\text{Green}} + \rho\_{\text{NIR}}}$$
\begin{itemize}
  \item \textbf{Sentinel-2 Band Map}: $\frac{B3 - B8}{B3 + B8}$
  \item \textbf{Landsat-8/9 Band Map}: $\frac{B3 - B5}{B3 + B5}$
  \item \textbf{Reflectance Range}: $[-1.0, 1.0]$. Open water features (such as tailing dams and lithium brine evaporation ponds) display positive NDWI values ($>0.1$), whereas soil and vegetation show negative values.

\end{itemize}
\subsection{Bare Soil Index (BSI)}
BSI integrates SWIR, red, NIR, and blue bands to isolate exposed soil and rock surfaces, providing an indicator for deforestation, mining activity, and soil disturbance:
$$\text{BSI} = \frac{(\rho\_{\text{SWIR1}} + \rho\_{\text{Red}}) - (\rho\_{\text{NIR}} + \rho\_{\text{Blue}})}{(\rho\_{\text{SWIR1}} + \rho\_{\text{Red}}) + (\rho\_{\text{NIR}} + \rho\_{\text{Blue}})}$$
\begin{itemize}
  \item \textbf{Sentinel-2 Band Map}: $\frac{(B11 + B4) - (B8 + B2)}{(B11 + B4) + (B8 + B2)}$
  \item \textbf{Landsat-8/9 Band Map}: $\frac{(B6 + B4) - (B5 + B2)}{(B6 + B4) + (B5 + B2)}$
  \item \textbf{Reflectance Range}: $[-1.0, 1.0]$. Exposed mining soils, quarries, and agricultural fields show high positive BSI values ($>0.15$), contrasting with healthy vegetation ($<0.0$).

\end{itemize}
\subsection{VIIRS Night-time Lights (NTL) Radiance Calibration}
The Suomi-NPP VIIRS Day/Night Band (DNB) captures nocturnal light emissions. Unlike legacy sensors, the VIIRS DNB undergoes absolute radiometric calibration, reporting radiance in units of nanowatts per square centimeter per steradian:
$$\text{Radiance} = \text{Pixel Value} \times 10^{-9} \text{ W}\cdot\text{cm}^{-2}\cdot\text{sr}^{-1}$$
Nocturnal radiance values correlate directly with economic factors such as electricity consumption, traffic volume, and industrial output.

---

\section{Unsupervised Anomaly Detection Mathematics}

\subsection{Spatial Anomaly Detection (Isolation Forest)}
For multi-spectral pixels in mining and industrial Regions of Interest (ROIs), an unsupervised \textbf{Isolation Forest} is trained. The core premise is that anomalies are few and spectrally distinct, making them easier to isolate than normal pixels.

Let $X = \{x\_1, x\_2, \dots, x\_N\}$ be a dataset of $N$ spatial observations in a $D$-dimensional space, where $D = [\text{NDVI}, \text{NDWI}, \text{BSI}, \rho\_{\text{SWIR1}}]$. The algorithm recursively splits the dataset by randomly selecting a feature and a random split value between the minimum and maximum values of that feature, constructing an ensemble of Isolation Trees ($iTrees$).

The anomaly score $s(x, n)$ for an observation $x$ is defined as:
$$s(x, n) = 2^{-\frac{\mathbb{E}(h(x))}{c(n)}}$$
Where:
\begin{itemize}
  \item $h(x)$ is the path length (number of edges traversed from the root node to a terminating leaf node) of sample $x$ in an individual tree.
  \item $\mathbb{E}(h(x))$ is the expected (average) path length of $x$ across the forest of $T$ trees.
  \item $c(n)$ is the average path length of an unsuccessful search in a Binary Search Tree (BST) built from $n$ nodes, serving as a normalization factor:
\end{itemize}
$$c(n) = 2\ln(n - 1) + 2\gamma - \frac{2(n - 1)}{n}$$
Here, $\gamma \approx 0.5772156649$ is the \textbf{Euler-Mascheroni constant}.

\textbf{Anomaly Thresholding}:
\begin{itemize}
  \item If $s(x, n) \approx 1.0$: The sample exhibits short path lengths across the forest, indicating that it is easily isolated and flagged as an anomaly.
  \item If $s(x, n) < 0.5$: The sample has long path lengths, indicating it falls within the normal land cover profile.
  \item The contamination parameter is set to $\alpha = 0.08$, defining the proportion of expected outliers.

\end{itemize}
\subsection{Temporal Anomaly Detection (Z-Score)}
To detect economic and industrial shocks in night-time light radiance time series, a rolling statistical Z-score is calculated:
$$Z\_t = \frac{Y\_t - \mu\_{(t-W, t-1)}}{\sigma\_{(t-W, t-1)}}$$
Where:
\begin{itemize}
  \item $Y\_t$ is the night-time lights radiance at month $t$.
  \item $\mu\_{(t-W, t-1)}$ is the historical mean calculated over a rolling window $W$ (configured as $W = 12$ months to account for annual cycles).
  \item $\sigma\_{(t-W, t-1)}$ is the standard deviation over the same rolling window $W$.

\end{itemize}
\textbf{Threshold Rules}:
\begin{itemize}
  \item $Z\_t > 2.5$: Represents a significant positive anomaly, suggesting rapid industrial expansion or infrastructure development.
  \item $Z\_t < -2.5$: Indicates a significant negative anomaly, suggesting factory closures, power outages, or economic contractions.

\end{itemize}
---

\section{Statistical Rigor \& Validation Matrix}
To validate the reliability of detected satellite anomalies, three statistical metrics are calculated:

\subsection{Spatial Classification Performance}
Using verified ground-truth shapes (e.g., official mining registry maps), the spatial anomalies are evaluated using a confusion matrix:
$$\text{Precision} = \frac{\text{True Positives (TP)}}{\text{True Positives (TP)} + \text{False Positives (FP)}}$$
$$\text{Recall} = \frac{\text{True Positives (TP)}}{\text{True Positives (TP)} + \text{False Negatives (FN)}}$$
$$\text{F1-Score} = 2 \cdot \frac{\text{Precision} \cdot \text{Recall}}{\text{Precision} + \text{Recall}}$$

\subsection{Temporal Economic Correlation}
To evaluate the correlation between night-time lights radiance anomalies ($Z\_t$) and official quarterly GDP growth, the pipeline computes Pearson's product-moment correlation coefficient ($r$):
$$r = \frac{\sum\_{i=1}^n (X\_i - \bar{X})(Y\_i - \bar{Y})}{\sqrt{\sum\_{i=1}^n (X\_i - \bar{X})^2 \sum\_{i=1}^n (Y\_i - \bar{Y})^2}}$$
Where $X\_i$ is the satellite radiance anomaly score, $Y\_i$ is the GDP growth rate, and $n$ is the number of observation quarters. 

A two-tailed Student's t-test is used to calculate the $p$-value, ensuring the correlation is statistically significant ($p < 0.05$).

\subsection{Spatial Autocorrelation (Moran's I)}
To account for spatial dependency—the principle that neighboring pixels are more likely to exhibit similar anomaly scores—we compute Moran's $I$ coefficient to verify spatial clustering:
$$I = \frac{N}{S\_0} \frac{\sum\_{i=1}^N \sum\_{j=1}^N w\_{ij}(x\_i - \bar{x})(x\_j - \bar{x})}{\sum\_{i=1}^N (x\_i - \bar{x})^2}$$
Where $N$ is the number of spatial units indexed by $i$ and $j$, $x$ is the anomaly score, $\bar{x}$ is the mean, $w\_{ij}$ is the spatial weight matrix, and $S\_0$ is the sum of all weights:
$$S\_0 = \sum\_{i=1}^N \sum\_{j=1}^N w\_{ij}$$
Moran's $I > 0$ indicates spatial clustering, while $I < 0$ indicates spatial dispersion, helping validate the structural consistency of detected mining footprints.

---

\section{Reproducibility Protocol}
To guarantee that third-party researchers can replicate all findings, the methodology implements:
\begin{enumerate}
  \item \textbf{Software Environment Locks}: All Python packages and system libraries are locked in a Docker configuration file (\href{file:///d:/CODING/tesis/pipeline/Dockerfile}{Dockerfile}) and \href{file:///d:/CODING/tesis/pipeline/requirements.txt}{requirements.txt}.
  \item \textbf{Algorithmic Seeds}: Random number generation seeds in scikit-learn's Isolation Forest are locked (\texttt{random\_state=42}), ensuring deterministic decision-tree splits.
  \item \textbf{Data Archives}: Raw coordinates, sensor metadata, and processing parameters are saved in a static \href{file:///d:/CODING/tesis/docs/data/regions.geojson}{regions.geojson} file, allowing identical geographic bounding boxes to be queried in future runs.

\end{enumerate}